\section{Notation, easy directions, and Corollaries~A and~B}\label{sec:notation}

\subsection{Notation}\label{sec:notation:basic}

All graphs are finite, simple, and undirected. Let $\graph$ be a graph
of order $n$, write $|E(\graph)|$ for its edge count, and let
$\Adj(\graph)$ denote its (symmetric, $0/1$) adjacency matrix. We
order the eigenvalues of $\Adj(\graph)$ decreasingly,
$\lambda_{1}(\graph) \ge \cdots \ge \lambda_{n}(\graph)$, and define
\[
\splus(\graph) := \sum_{i: \lambda_{i}(\graph) > 0} \lambda_{i}(\graph)^{2},
\qquad
\sminus(\graph) := \sum_{i: \lambda_{i}(\graph) < 0} \lambda_{i}(\graph)^{2}.
\]
The trace identity
$\tr\bigl(\Adj(\graph)^{2}\bigr) = \splus(\graph) + \sminus(\graph) =
2|E(\graph)|$ is standard. The \emph{positive} and \emph{negative
inertia} are
$n^{+}(\graph) := |\{i : \lambda_{i}(\graph) > 0\}|$ and
$n^{-}(\graph) := |\{i : \lambda_{i}(\graph) < 0\}|$.

A \emph{$2$-tree} is a graph constructed inductively starting from
$K_{2}$ by repeatedly adding a new vertex adjacent to both endpoints of
an existing edge. Equivalently, a $2$-tree is a maximal chordal graph
of treewidth $2$; every maximal clique is a triangle. The
\emph{clique tree} $\cliqtree(\graph)$ of a $2$-tree has the
$n - 2$ triangles of $\graph$ as nodes and an edge between any two
triangles sharing an edge of $\graph$; $\cliqtree(\graph)$ is itself a
tree. A vertex $v$ of $\graph$ is \emph{simplicial} if its neighbourhood
$N_{\graph}(v)$ is a clique. The simplicial degree-$2$ vertices of a
$2$-tree are precisely the apex vertices of leaf triangles of
$\cliqtree(\graph)$; for such a vertex $v$ with neighbours $\{a, b\}$,
the edge $\{a, b\}$ (which lies in $\graph$) is its \emph{supporting
edge}. We write $\Hgraph := \graph - v$ for the deletion of $v$, and
note that $\Hgraph$ is itself a $2$-tree on $n - 1$ vertices.

\subsection{Easy forward implications}\label{sec:notation:easy}

The forward direction of \Cref{conj:92} is elementary. If $\graph = T$
is a tree on $n$ vertices, $T$ is bipartite, so $\spec(\Adj(T))$ is
symmetric about $0$ and the trace identity gives
$\splus(T) = \sminus(T) = |E(T)| = n - 1$. If $\graph = K_{n}$, the
adjacency matrix is $J - I$ where $J$ is the all-ones matrix, with
spectrum $\{n - 1\} \cup \{-1\}^{(n-1)}$; hence
$\splus(K_{n}) = (n - 1)^{2}$ and $\sminus(K_{n}) = n - 1$. The
substantive content of \Cref{conj:92} is the converse direction in each
part.

\subsection{Corollary~A: Conjecture~9.2(i) for connected claw-free graphs}
\label{sec:notation:corA}

\begin{corollary}[Conjecture~9.2(i) for connected claw-free graphs]
\label{cor:A}
If $\graph$ is a connected claw-free graph with $\splus(\graph) = n - 1$,
then $\graph$ is a tree.
\end{corollary}

\begin{proof}
We split on the maximum degree $\Delta(\graph)$.

If $\Delta(\graph) \ge 3$, Theorem~1.1 of \cite{AKMPZ2025} gives
$\splus(\graph) \ge n$, contradicting the hypothesis
$\splus(\graph) = n - 1$.

If $\Delta(\graph) \le 2$, $\graph$ is a path or a cycle. Paths are
trees, and the conclusion holds. For a cycle $C_{m}$ on $m \ge 3$
vertices, the spectrum of $\Adj(C_{m})$ is
$\{2 \cos(2\pi k / m) : 0 \le k < m\}$, with $\lambda_{1}(C_{m}) = 2$
and the remaining eigenvalues distributed symmetrically around $0$
when $m$ is even. A short calculation using
$\sum_{k} \cos^{2}(2\pi k / m) = m/2$ for $m \ge 3$ gives
$\splus(C_{m}) \ge m \ne m - 1$, contradicting the hypothesis. Hence
$\graph$ must be a path, i.e.\ a tree.
\end{proof}

\begin{remark}\label{rem:Pjkl}
The proof above delegates the heavy lifting on claw-free graphs of
maximum degree $\ge 3$ — in particular the infinite family
$P(j, k, \ell)$ of subdivided-$K_{4}$-like 2-trees, where claw-freeness
coexists with unicyclic structure — to Theorem~1.1 of \cite{AKMPZ2025}.
We do not reprove that case here.
\end{remark}

\subsection{Corollary~B: Conjecture~9.2(i) for diameter $\le 2$}
\label{sec:notation:corB}

\begin{corollary}[Conjecture~9.2(i) for $\diam(\graph) \le 2$]
\label{cor:B}
If $\graph$ is a connected graph with $\diam(\graph) \le 2$ and
$\splus(\graph) = n - 1$, then $\graph$ is a tree.
\end{corollary}

\begin{proof}
We split on the diameter. If $\diam(\graph) = 0$ then $\graph = K_{1}$,
trivially a tree. If $\diam(\graph) = 1$ then $\graph = K_{n}$, so
$\splus(\graph) = (n - 1)^{2}$; the hypothesis
$\splus(\graph) = n - 1$ then forces $n = 2$, and $\graph = K_{2}$ is a
tree.

If $\diam(\graph) = 2$, Theorem~1.2 of \cite{AKMPZ2025} gives
$\splus(\graph) \ge n$ except for $\graph \in \{K_{1, n-1}, C_{5}\}$.
The hypothesis $\splus(\graph) = n - 1$ therefore excludes everything
outside these two cases. The star $K_{1, n-1}$ is a tree, with
$\splus(K_{1, n-1}) = n - 1$ — consistent with the hypothesis. The
5-cycle $C_{5}$ has spectrum
$\{2, 2\cos(2\pi/5), 2\cos(2\pi/5), 2\cos(4\pi/5), 2\cos(4\pi/5)\}$,
giving $\splus(C_{5}) = 4 + 2(2\cos(2\pi/5))^{2} = 4 + (3 - \sqrt{5})
= 7 - \sqrt{5} \approx 4.764 > 4 = n - 1$, contradicting the hypothesis. Hence
the only graph compatible with $\splus = n - 1$ at diameter~$2$ is the
star, which is a tree.
\end{proof}

\subsection{The slack: why these corollaries are not the paper}
\label{sec:notation:slack}

Corollaries~\ref{cor:A} and~\ref{cor:B} are short consequences of
Theorems~1.1 and~1.2 of \cite{AKMPZ2025}: each rules out
$\splus = n - 1$ on its class by a strict inequality already proved in
the source paper, and then exhausts the exceptional cases by hand.
These deductions verify the toolkit and serve as an entry point to the
harder content of \Cref{sec:lprime,sec:subfamily,sec:moment-form}; they
do not by themselves constitute new results.
